\chapter{Approfondimento sul CAD service}

\section{Ruolo del servizio}

Il \texttt{cadService} e' il cuore applicativo di KompozeR. Non si limita a salvare dati di configurazione, ma interpreta il progetto dello scaffale come una sequenza di vincoli e trasformazioni. Il suo compito e' mantenere una \texttt{Configuration} valida sia dal punto di vista strutturale sia dal punto di vista commerciale, fino alla generazione della distinta base finale.

In pratica il servizio prende una configurazione che evolve a step e controlla che ogni passaggio rispetti le regole del catalogo e la geometria della composizione. Quando la configurazione e' completa, il CAD deriva la BOM e la passa al carrello.

\section{Modello di dominio}

Il modello reale del servizio ruota attorno all'aggregate root \texttt{Configuration}, che contiene:

\begin{itemize}
\item identita' e proprietario;
\item nome e stato del progetto;
\item categoria selezionata;
\item vincoli di ambiente;
\item piano colonne;
\item disegni colonna per colonna;
\item componenti derivati in BOM;
\item versione ottimistica della configurazione.
\end{itemize}

L'implementazione valida anche gli invarianti di base: identita' non vuota, proprietario presente, versione almeno pari a 1, coerenza tra numero colonne e piano colonne, livelli crescente e spessori positivi.

\section{Flusso step-based}

Il percorso reale nel codice e' il seguente:

\begin{enumerate}
\item creazione di una configurazione in stato \texttt{DRAFT};
\item definizione dell'ambiente fisico massimo/minimo;
\item selezione della categoria \texttt{TONDO}, \texttt{QUADRO} o \texttt{KUBE};
\item definizione del piano colonne con larghezze ammesse dal catalogo;
\item aggiornamento del design colonna per colonna;
\item calcolo delle opzioni successive disponibili per ciascuna colonna;
\item derivazione della BOM quando il design e' completo;
\item finalizzazione e sincronizzazione con il carrello.
\end{enumerate}

Questo flusso e' implementato dai use case dedicati: \texttt{CreateConfiguration}, \texttt{SetEnvironment}, \texttt{SetCategory}, \texttt{SetColumnPlan}, \texttt{UpdateDesign}, \texttt{FinalizeConfiguration}, \texttt{ReorderConfiguration}, \texttt{GetConfiguration}, \texttt{ListConfigurations} e \texttt{ListNextOptions}.

\section{Logica geometrica}

Il CAD non accetta modifiche arbitrarie. La logica geometrica e' affidata a \texttt{SpineModel} e \texttt{deriveBom}. Le regole principali sono:

\begin{itemize}
\item la shelf thickness e' server-driven e fissata a 20 mm;
\item i livelli di una colonna devono crescere in modo strettamente monotono;
\item le colonne adiacenti non possono condividere livelli incompatibili;
\item i segmenti tra livelli successivi devono corrispondere a montanti disponibili nel catalogo;
\item il terminale superiore deve entrare nel vincolo di altezza massima.
\end{itemize}

La funzione \texttt{validateColumnCandidate} valuta la proposta di aggiungere un nuovo livello in una colonna e controlla sia il conflitto di adiacenza sia la validita' delle spine coinvolte. La funzione \texttt{validateColumnDesigns} verifica l'intero snapshot delle colonne, mentre \texttt{deriveSpineBom} e \texttt{deriveBom} trasformano il design valido in componenti commerciali.

\section{Interazione con il catalogo}

Il CAD consulta il catalogo attraverso \texttt{HttpCatalogRulesProvider}. Questa integrazione traduce i prodotti catalogati in regole utili al configuratore:

\begin{itemize}
\item larghezze dei ripiani;
\item altezze disponibili per piedini, montanti e terminali;
\item SKU e nomi dei componenti;
\item prezzo associato a ciascun componente.
\end{itemize}

Il servizio usa queste regole per determinare se una colonna puo' essere costruita e per calcolare la distinta finale. In questo senso il CAD dipende dal catalogo per il \emph{vocabolario fisico} del prodotto, ma conserva la propria logica decisionale.

\section{Derivazione della BOM}

La BOM e' costruita come insieme aggregato di componenti commerciali. La funzione \texttt{deriveBom} procede in tre fasi:

\begin{itemize}
\item conta i ripiani per colonna e larghezza;
\item calcola i componenti di spine condivise, cioe' piedini, terminali e montanti;
\item aggrega gli item con lo stesso SKU prima di restituire il risultato.
\end{itemize}

Questo passaggio e' centrale per il resto dell'applicazione: il cart service riceve la BOM gia' pronta, il frontend mostra il riepilogo prezzo, e la finalizzazione puo' essere bloccata se la distinta non e' derivabile.

\section{Stati della configurazione}

Nel codice reale gli stati sono:

\begin{itemize}
\item \texttt{DRAFT};
\item \texttt{ENVIRONMENT\_DEFINED};
\item \texttt{CATEGORY_SELECTED};
\item \texttt{COLUMNS_DEFINED};
\item \texttt{DESIGN_IN_PROGRESS};
\item \texttt{READY_FOR_FINALIZE};
\item \texttt{FINALIZED}.
\end{itemize}

Il flusso non e' solo descrittivo: ogni use case controlla lo stato corrente e impedisce transizioni non ammesse. Ad esempio non si puo' cambiare ambiente o categoria quando la configurazione e' finalizzata, e non si puo' finalizzare senza componenti derivabili.

\section{Perche' il CAD e' il core}

Il CAD e' il core per tre motivi:

\begin{itemize}
\item concentra la logica di business piu' specifica e piu' complessa;
\item governa il passaggio dall'idea del cliente a una distinta realmente acquistabile;
\item influenza direttamente carrello, notifiche, chatbot e reporting.
\end{itemize}

Se il CAD sbaglia la geometria o la BOM, il resto del sistema lavora su dati non coerenti. Per questo nel progetto attuale il servizio e' stato modellato con maggiore attenzione rispetto agli altri microservizi, e questo documento lo tratta come il centro del dominio applicativo.
